\documentclass{misisarticle}
\usepackage{flafter}
\usepackage{tikz}
\usetikzlibrary{arrows.meta,positioning,shapes.geometric}
\addto\captionsrussian{\renewcommand{\refname}{ЛИТЕРАТУРА}}

\grnti{20.23.21}
\asjc{1706}
\oecd{1.2}

\title{Поиск аналогов косметических средств на основе анализа состава и семантических характеристик продукции}
\misisauthor{М.~А.~Меркулов}{кафедра инженерной кибернетики НИТУ МИСИС}{Москва, Россия}{m2306898@edu.misis.ru}
\misisauthor{В.~В.~Семенов}{кафедра инженерной кибернетики НИТУ МИСИС}{Москва, Россия}{m2306898@edu.misis.ru}

\abstract{Рассматривается задача поиска аналогов косметических средств с учетом их ингредиентного состава и заявленных текстовых характеристик. Цель в повышении релевантности поисковой выдачи для различных типов пользовательских запросов, включая INCI-составы и запросы на естественном языке. Предложен гибридный метод поиска, в котором запросы по составу обрабатываются с помощью BM25, запросы на естественном языке с использованием семантического поиска на базе BGE-M3, а для повышения качества ответов в более тяжелых случаях применяется LLM-реранжирование найденных кандидатов. Используется самостоятельно собранный датасет, который содержит сведения о составе, описании и характеристиках косметической продукции. Эффективность рассмотренных методов оценивается по метрике NDCG@10 и времени обработки пользовательского запроса.}
\keywords{поиск аналогов, INCI, BM25, BGE-M3, LLM}

\begin{document}
\maketitle

\section{Введение}

В работе рассматривается задача подбора наиболее релевантных косметических средств, исходя из пользовательского запроса: по составу, по назначению и по заявленным характеристикам - такой запрос может содержать список ингредиентов исходного продукта, название средства, описание желаемых свойств или указание на продукт, для которого требуется найти аналог. Например, запрос "увлажняющий крем для жирной кожи" задает назначение и тип кожи, а список из ингредиентов "вода, глицерин, ниацинамид, пантенол" непосредственно задает ингредиенты для сравнения.

В зависимости от того, как пользователь формулирует запрос - меняются и критерии поиска. При сопоставлении составов, основное внимание уделяется совпадению названий компонентов. При обработке свободного описания, необходимо определить насколько содержание запроса соответствует описанию товара. Кроме того, поиск конкретного косметического средства отличается от подбора похожего продукта. В первом случае бренд может быть важным критерием, а во втором большее значение имеют характеристики и свойства косметического средства.

В существующих исследованиях в области косметических рекомендаций, рассматривают извлечение характеристик из ингредиентов~\cite{liu2024}, подбор средств с учетом пользовательских данных~\cite{rajegowda2024} и применение экспертных правил~\cite{fedorov2025}. В отличии от этих направлений, данная работа сосредоточена на организации поиска в каталоге при неоднородных запросах. Основная гипотеза состоит в том, что выбор поискового метода с учетом формы запроса позволяет получить более релевантную результаты, чем применение одного универсального способа.

Основным вкладом работы заключается в разработке и реализации схемы маршрутизации между поиском по ингредиентам и семантическим поиском, а также дополнительного этапа переранжирования с критериями для разных типов запроса. Для экспериментальной проверки этого подхода, реализован программный прототип и набор тестовых запросов. Предлагается оценивать качество ранжирования совместно с задержкой обработки, включая затраты на подключаемые языковые модели, что даст более объективную оценку.

\section{Анализ связанных работ} 

\subsection{Поиск и рекомендации косметической продукции}

S.~Liu, R.~Suresh и A.~Banitalebi-Dehkordi~\cite{liu2024} предлагают определять характеристики косметических товаров по их составу и названий. Авторы используют обучаемую модель, а полученные характеристики рассматривают как основу для рекомендаций и объяснения выбора товаров. Эта работа показывает способ перехода от состава к признакам, которые можно использовать при рекомендациях.

G.~M.~Rajegowda и соавторы~\cite{rajegowda2024} описывают систему рекомендаций по уходу за кожей лица. Пользователь заполняет анкету, а сверточная нейронная модель анализирует изображение его лица. Полученные данные используются для подбора средств с учетом характеристик кожи. В работе косметичиеские средства подбираются по текстовому запросу. При оценке качества поиска проверяется, насколько найденные товары соответствуют запросу и занимают ли наиболее подходящие из них первые позиции в выдаче. 

А.~А.~Федоров~\cite{fedorov2025} рассматривает систему подбора уходовых средств с учетом типа кожи, возраста, бюджета и других характеристик пользователя. Для настоящей задачи этот подход представляет интерес как пример явного задания критериев подбора. В разрабатываемой системе критерии используются при выборе поискового пути и при дополнительном ранжировании кандидатов.

\subsection{Работы из сборника НИТУ МИСИС}

Д.~В.~Береснев~\cite{beresnev2023} сравнивает применение предобученных моделей при разработке корпоративного чат-бота. В частности, рассматривается семантический поиск с использованием векторных представлений и ближайших соседей. Для настоящего исследования полезен подход к раздельному выбору моделей для различных задач обработки текста.

Е.~С.~Четвертков~\cite{chetvertkov2026} сравнивает мультимодальные архитектуры для поиска трехмерных моделей LEGO по текстовым запросам. Общим с настоящей работой является исследование готовых моделей в специализированной предметной области и сопоставление качества с вычислительными затратами. При этом данные и способы представления объектов в этих задачах различаются.

Рассмотренные исследования задают методы извлечения признаков, семантического поиска и переранжирования. В данной работе вопрос состоит в том, насколько полезно выбирать между поиском по ингредиентам и поиском по тексту для разных пользовательских формулировок в одном каталоге косметической продукции.

\section{Постановка задачи}

Пусть $D=\{d_1,\ldots,d_N\}$ - каталог товаров. Каждому товару $d$ соответствуют идентификатор, название, бренд, описание, характеристики и состав, если эти поля доступны. Пользователь задает текстовый запрос $q$. Требуется сформировать упорядоченный список $R_k(q)\subseteq D$ из не более чем $k$ товаров, соответствующих запросу.

В работе под аналогом понимается косметические средства, близкие к указанному в запросе составу или соответствующее описанным требованиям. Для запроса по названию допускается поиск самого продукта. Для запроса на замену требуется оценивать сходства с исходным продуктом. Поэтому качество поиска в этих двух случаях следует оценивать раздельно. При этом совпадение текстовых признаков не означает, что средства имеют одинаковую рецептуру или одинаковое действие. 

Для каждого товара строятся два представления: список нормализованных ингредиентов $I(d)$ и текст $T(d)$, объединяющий название, тип продукта и описание. Функция $h(q)$ определяет, является ли запрос списком ингредиентов. Базовое ранжирование имеет вид
\begin{equation}
R_k^{(0)}(q)=
\begin{cases}
\operatorname{Top}_k B(q,I(d)), & h(q)=1,\\
\operatorname{Top}_k S(q,T(d)), & h(q)=0,
\end{cases}
\label{eq:routing}
\end{equation}
где $B$ - оценка BM25, $S$ - близость плотных векторов. Обозначение $\operatorname{Top}_k$ означает выбор товаров с наибольшей оценкой и их упорядочение по убыванию этой оценки.

Для расширенного варианта сначала отбирается не более 20 кандидатов по формуле~\eqref{eq:routing}, затем языковая модель уточняет их порядок. Она не добавляет товары за пределами найденного набора. Поэтому качество первого этапа ограничивает возможности последующего переранжирования.

Для проверки гипотези все методы сравниваются на одном и том же каталоге товаров и наборе запросов. Сравниваются BM25 по ингредиентам, семантический поиск по текстам карточек товаров и маршрутизация с дополнительным LLM-переранжированием. Качество выдачи оценивается по метрике NDCG@10, а скорость работы по времени обработки запроса. Результаты анализируются как по всему набору запросов, так и отдельно для каждого типа.

\section{Данные и их подготовка}

Используется самостоятельно собранный каталог косметической продукции. Источниками данных служат интернет-магазины "Золотое Яблоко", "РИВ ГОШ", "ЛЭТУАЛЬ" и "hollyshop". Записи содержат идентификатор, URL, название, бренд, цену на момент сбора, описание, способ применения, характеристики, страну происхождения, производителя и состав.

Подготовка включает извлечение структурированных характеристик из строки вида "ключ: значение", очистку состава и получение списка ингредиентов. Для состава удаляются кириллические пояснения в круглых и квадратных скобках, унифицируются поддерживаемые разделители, регистр и пробелы. Затем выполняется разбиение на компоненты. Например, "Grape seed oil" представляется единым лексическим токеном \texttt{grape\_seed\_oil}.

В данный момент дополнительно предусмотрена дополнительная обработка данных о товарах. Активные компоненты выделяются с помощью заранее заданного словаря. LLM извлекает из текстов карточек сведения о заявленных свойствах средств, типах кожи и задачах ухода и записывает их в виде меток. Полученные метки сверяются со списком допустимых значений и приводятся к единому виду по таблицам соответствий. 

В базовой маршрутизации используется состав и текстовое представление карточки. Дополнительные метки не используются непосредственно в ее формуле оценки. При этом система не переводит русские названия ингредиентов в обозначения INCI. Из-за этого она может распознать русскоязычный запрос как список ингредиентов, но не найти совпадений с теми же компонентами, которые записаны в каталоге по INCI-составу.

\section{Описание метода}

\subsection{Маршрутизация запросов}

Система распознает список ингредиентов по заранее заданным правилам. Она проверяет наличие запятых, число частей списка и вхождения терминов из фиксированного русско-английского словаря. Для выбора INCI-пути требуется не менее четырех частей и не менее двух словарных вхождений либо не менее шести частей и хотя бы одно такое вхождение. Во всех остальных случаях используется семантический поиск.

Таким образом, короткий запрос "крем с ниацинамидом" обрабатывается как естественный язык. Система выбирает метод поиска по заданным правилам без обращения к LLM. В данной работе этот выбор выполняется до обработки состава: если выбран поиск по ингредиентам, система затем очищает состав и приводит названия компонентов к единому формату. Схема показана на рис.~\ref{fig:pipeline}.

\begin{figure}[!htbp]
\centering
\begin{tikzpicture}[
  >=Latex,
  box/.style={draw,rounded corners=2pt,align=center,font=\footnotesize,minimum height=0.7cm,text width=2.9cm},
  smallbox/.style={box,text width=2.9cm},
  line/.style={->,line width=0.45pt},
  every edge quotes/.style={font=\scriptsize}
]
\node[box] (q) at (0,0) {Запрос пользователя};
\node[box] (router) at (0,-1.05) {Эвристический детектор INCI};
\node[smallbox] (bm) at (-1.8,-2.25) {BM25 по ингредиентам};
\node[smallbox] (dense) at (1.8,-2.25) {BGE-M3 по тексту карточки};
\node[box] (cand) at (0,-3.45) {Ранжированный набор кандидатов};
\node[box] (llm) at (0,-4.55) {Опциональное\newline переранжирование LLM};
\node[box] (out) at (0,-5.65) {Итоговая выдача};
\draw[line] (q)--(router);
\draw[line] (router.south west)--node[left,font=\scriptsize]{список} (bm.north);
\draw[line] (router.south east)--node[right,font=\scriptsize]{текст} (dense.north);
\draw[line] (bm.south)--(cand.north west);
\draw[line] (dense.south)--(cand.north east);
\draw[line] (cand)--(llm);
\draw[line] (llm)--(out);
\draw[line] (cand.west)--++(-2.1,0)|-node[pos=0.2,left,font=\scriptsize]{без LLM} (out.west);
\end{tikzpicture}
\caption{Маршрутизация и дополнительное переранжирование кандидатов.}
\label{fig:pipeline}
\end{figure}

\subsection{Лексический поиск по ингредиентам}

Состав каждого товара преобразуется в последовательность токенов, каждый из которых соответствует одному компоненту. На основе этой коллекции строится BM25-индекс. Если система выбирает поиск по составу, запрос обратывается по тем же правилам. Затем определяется оценка товара путем совпадения токенов с учетом частоты компонентов и длины списка \cite{robertson2009}.

Для расчета используется \texttt{BM25Okapi} из библиотеки \texttt{rank\_bm25}. В выдачу включаются только твоары с оценкой выше нуля. Обычный BM25 рассматривает частоты токенов и не вводит отдельный вес позиции компонента в составе. Поэтому он не определяет концентрации компонентов и не учитвает порядок их перечисления при сравнеии составов

\subsection{Семантический поиск}

Текст товара строится объединением названия, типа продукта и описания. В данный момент описание ограничивается первыми 1500 символами. С помощью модели BGE-M3 для каждого текста заранее вычисляется нормализованный вектор ~\cite{chen2024}. При поступлении запроса той же моделью вычисляется его нормализованное представление.

Пусть $e_q$ и $e_d$ - векторы запроса и товара. При $\|e_q\|_2=\|e_d\|_2=1$ оценка косинусной близости равна скалярному произведению:
\begin{equation}
S(q,T(d))=e_d^{\mathsf T}e_q.
\label{eq:cos}
\end{equation}
Для матрицы векторов каталога $X$ все оценки получаются операцией $Xe_q$. Затем система выбирает товары с наибольшими оценками. В данный момент запрос сравнивается со всем каталогом. Приближенный поиск ближайших соседей не используется.

\subsection{Типизированное LLM-переранжирование}

После основного поиска порядок найденных товаров может быть дополнительно уточнен с помощью LLM. Для этого дополнительного этапа LLM получает не более 20 кандидатов базового поиска. Для каждого передается идентификатор, название, бренд, тип продукта, первые шесть ингредиентов и сокращенное описание. Модель получает весь список и возвращает JSON-массив идентификаторов в новом порядке ~\cite{sun2023}.

Критерии определяются типом запроса. Для поиска по ингредиенту приоритетом является упомянутый компонент и категория товара. Для поиска по названию - соответствие названию и указанному бренду. Для поиска аналога - класс продукта и заявленные свойства. Для описания потребности - соответствие описанию и типу кожи. Эти критерии задаются модели в виде инструкций для упорядочения товаров. Они не служат отдельной првоеркой достоверности сведений в карточках.

В данный момент тип запроса передается отдельной структурой, полученной модулем анализа запроса. Для типов INCI и "категория" переранжирование пропускается. При некорректном ответе или ошибке API сохраняется исходный порядок. Идентификаторы за пределами набора кандидатов и повторы отбрасываются, а пропущенные моделью кандидаты дополняют список.

\section{Протокол оценки}

Для оценки был подготовлен набор из 130 запросов шести типов (табл.~\ref{tab:queries}). Запросы сгенерированы LLM на основе карточек каталога и имитируют различные пользовательские сценарии. Для получения разметки объединяются кандидаты нескольких поисковых методов: пары "запрос - товар" оцениваются языковой моделью по шкале 0, 1, 2. 

\begin{table}[!htbp]
\caption{Состав доступного набора запросов}
\label{tab:queries}
\centering
\begin{tabular}{@{}p{0.76\columnwidth}r@{}}
\toprule
Тип запроса & Число \\
\midrule
Список ингредиентов & 25 \\
Название продукта & 25 \\
Категория товара & 20 \\
Категория с указанием ингредиента & 25 \\
Описание потребности & 20 \\
Поиск аналога продукта-референса & 15 \\
\midrule
Всего & 130 \\
\bottomrule
\end{tabular}
\end{table}

Основная метрика - NDCG@10. Она учитывает соответствие первых десяти товаров запросу и их положение в выдаче. Как и в программной реализации, оценки ревалентности используются в формуле напрямую, без дополнительного преобразования.
\begin{equation}
\begin{aligned}
\mathrm{DCG}@k&=\sum_{i=1}^{k}\frac{r_i}{\log_2(i+1)},\\
\mathrm{NDCG}@k&=\frac{\mathrm{DCG}@k}{\mathrm{IDCG}@k},
\end{aligned}
\label{eq:ndcg}
\end{equation}
где $r_i\in\{0,1,2\}$ - оценка товара на позиции $i$, а IDCG - значение для идеального порядка известных оценок. При нулевом IDCG значение метрики в коде принимается равным нулю. Неразмеченные кандидаты учитываются как имеющие нулевую релевантность. Итоговая метрика усредняется по запросам, дополнительно приводится разбивка по их типам.

Для измерения задержки фиксируется медиана и 95-й процентиль времени обработки. Скрипты получают структурированный тип запроса из заранее подготовленного кэша. Поэтому время вызова поискового метода следует отделять от полного времени обработки, включающего анализ запроса, обращение к внешней LLM и формирование ответа API. 

\begin{thebibliography}{9}

\bibitem{liu2024}
Liu S., Suresh R., Banitalebi-Dehkordi A.
Beauty Beyond Words: Explainable Beauty Product Recommendations Using Ingredient-Based Product Attributes // Proceedings of SURE 2024. CEUR Workshop Proceedings. 2024. Vol.~3873.
URL: \url{https://ceur-ws.org/Vol-3873/87.pdf}.

\bibitem{rajegowda2024}
Rajegowda G.~M., Spyridis Y., Villarini B., Argyriou V.
An AI-Assisted Skincare Routine Recommendation System in XR. 2024. arXiv:2403.13466.
URL: \url{https://arxiv.org/abs/2403.13466}.

\bibitem{fedorov2025}
Федоров А.~А. Разработка экспертной системы для персонализированного подбора уходовых косметических средств на основе индивидуальных характеристик пользователя // Новая наука в новом мире: сборник статей XI Международной научно-практической конференции. Петрозаводск: МЦНП «Новая наука», 2025. С.~184--188.
URL: \url{https://sciencen.org/assets/Kontent/Konferencii/Arhiv-konferencij/KOF-1289.pdf}.

\bibitem{robertson2009}
Robertson S., Zaragoza H. The Probabilistic Relevance Framework: BM25 and Beyond // Foundations and Trends in Information Retrieval. 2009. Vol.~3, no.~4. P.~333--389.
DOI: \url{https://doi.org/10.1561/1500000019}.

\bibitem{chen2024}
Chen J., Xiao S., Zhang P., Luo K., Lian D., Liu Z.
M3-Embedding: Multi-Linguality, Multi-Functionality, Multi-Granularity Text Embeddings Through Self-Knowledge Distillation // Findings of ACL 2024. 2024. P.~2318--2335.
DOI: \url{https://doi.org/10.18653/v1/2024.findings-acl.137}.

\bibitem{sun2023}
Sun W. et al. Is ChatGPT Good at Search? Investigating Large Language Models as Re-Ranking Agents // EMNLP 2023. 2023. P.~14918--14937.
DOI: \url{https://doi.org/10.18653/v1/2023.emnlp-main.923}.

\bibitem{beresnev2023}
Береснев Д.~В. Анализ подходов к использованию предобученных моделей в разработке корпоративных чат-ботов // Искусственный интеллект в промышленных, коммерческих, медицинских и финансовых приложениях: сборник статей научно-технического семинара студентов / под ред. А.~Р.~Ефимова. М.: НИТУ «МИСИС», 2023. Вып.~1. С.~68--77.
URL: \url{https://misis.ru/files/-/1a3a57deb59c39284cba56d0b54bb308/kik_2023.pdf}.

\bibitem{chetvertkov2026}
Четвертков Е.~С. Сравнительный анализ мультимодальных архитектур в задаче текстового поиска 3D-моделей LEGO // Искусственный интеллект в промышленных, коммерческих, медицинских и финансовых приложениях: сборник статей научно-технического семинара / под ред. Р.~Н.~Садекова. М.: НИТУ «МИСИС», 2026. Вып.~6. С.~141--146.
Авторская версия: \url{https://www.researchgate.net/publication/407039821}.
\end{thebibliography}
\end{document}